% Mục 6.2 (bản rút gọn). Rút từ sections/results_barrier.tex; nguồn từng số ghi ở đó
% (results/apple-m5/barrier.csv, ns mỗi episode = 1e9 / trung vị ops_per_s; scripts/report_numbers/rbc_analyze.py).
\subsubsection{RQ4: Chi phí một barrier episode}

\begin{rqanswer}\textbf{Trả lời RQ4.} Với chờ spin, một episode tốn thêm khoảng 53~ns cho mỗi thread
trong cụm 4 P-core. Độ dốc tăng mạnh khi vượt 4 thread (số P-core), gần gấp 6 theo trung vị. Với \PCname{sr}, số episode
mỗi giây giảm khoảng 2\,000 lần khi số thread vượt số core. Khi mỗi thread có core riêng, spin thuần
nhanh nhất hoặc ngang \PCname{sr_bo_yield} (200 so với 195~ns ở 4 thread). Khi vượt số core,
\PCname{sr_bo_yield} nhanh nhất. Số liệu của \PCname{apple-m5}. RQ6 kiểm lại trên hai máy còn
lại.\end{rqanswer}
% nguồn: barrier.csv, ns/episode p=4: sr 200,1, sr_bo_yield 195,2. sr: tốc độ p=10 / p=12 = 2\,067.

\PCfigure{apple-m5}{barrier_rate}

\emph{Bằng chứng.} Nhóm đo số episode mỗi giây khi giữa hai barrier không có việc gì
(Hình~\ref{fig:apple-m5:barrier_rate}). Từ 1 lên 2 thread, một episode của \PCname{sr} tăng từ 1,7~ns
lên 94,5~ns. Từ 2 tới 4 thread, mỗi thread thêm vào tốn khoảng 53~ns, từ 4 tới 10 thread khoảng 295~ns.
Từ 10 lên 12 thread, \PCname{sr} chậm đi khoảng 2\,070 lần, và mỗi thread vượt số core thêm khoảng 2~ms
vào một episode. Tới 4 thread, \PCname{sr}, \PCname{sr_block} và \PCname{sr_bo_yield} gần nhau,
\PCname{std_barrier} chậm hơn \PCname{sr} 1,6 lần, còn \PCname{sr_bo_spin}, \PCname{condvar} và
\PCname{sr_bo_sleep} chậm hơn 3,9--31 lần. Khi vượt số core, thứ tự đảo lại. Năm barrier không quay
thuần nhanh hơn \PCname{sr} khoảng 170--680 lần, và \PCname{sr_bo_yield} nhanh nhất.
% nguồn: barrier.csv, ns/episode sr p=1 1,70, p=2 94,5; p=4 200,1; p=10 1\,969; p=4 std_barrier 320,8 (1,60),
% sr_bo_spin 786,5 (3,93), sr_bo_sleep 6\,112 (30,5); p=12..20: 168..677 lần.

\emph{Vì sao.} Từ 2 thread, bộ đếm và cờ phải đi giữa các core ở mỗi episode. Phần tuyến tính đến từ
$p$ phép \texttt{fetch\_add} lần lượt trên cache line của bộ đếm và $p-1$ lần đọc lại cờ sau khi cờ đảo.
Chỗ gãy sau 4 thread trùng với số P-core, và kết quả lock cũng gãy ở cùng chỗ ngay cả ở ô đối chứng
(RQ2). Vì vậy nhóm cho rằng nó đến từ máy chứ không từ thuật toán barrier. Khi vượt số core, thread đã tới
barrier quay và giữ core tới hết lượt lập lịch, nên thread chưa tới phải chờ lượt sau mới được chạy, giống
MCS ở RQ2.

\PCname{sr_block} chờ bằng \texttt{std::atomic::wait}. Trong libc++ của Apple clang 21, lời gọi này đọc
tới khi đã chờ 4~$\mu$s rồi mới ngủ trong kernel. Tới 4 thread, episode ngắn hơn 4~$\mu$s nhiều, nên
\PCname{sr_block} gần bằng \PCname{sr}. Từ 6 thread nó chậm hẳn (16,1~$\mu$s ở 10 thread). Nhóm cho rằng
ngày càng nhiều episode có thread chờ quá 4~$\mu$s, ngủ, rồi phải được đánh thức qua kernel.
\PCname{sr_bo_yield} đọc cờ 64 lần rồi mới gọi \texttt{sched\_yield}. Khi vượt số core, thread chờ nhường
core, có thể cho thread chưa tới barrier, nên episode xong mà không phải chờ hết lượt lập lịch và không
phải trả một lần đánh thức qua kernel. Ở 12 thread, một episode của nó tốn 11,6~$\mu$s, so với 13,7~$\mu$s
của \PCname{std_barrier} và 18,8~$\mu$s của \PCname{sr_block}. \PCname{condvar} không có pha quay, nên
ngay ở 2 thread mỗi episode đã tốn 1,77~$\mu$s, và tới 4 thread mỗi thread thêm vào tốn khoảng
1,7~$\mu$s, cùng bậc với chi phí đánh thức 1,5~$\mu$s ước lượng ở RQ2.
% nguồn: header libc++ (atomic_sync.h); barrier.csv, ns/episode sr_block p=10 16\,064; p=12 sr_bo_yield 11555,
% std_barrier 13657, sr_block 18755; condvar p=2/4 1\,771/5\,175.

Backoff làm barrier chậm đi khi mỗi thread có core riêng (\PCname{sr_bo_spin} chậm hơn \PCname{sr}
2,3 tới 9,5 lần), vì thread chờ chỉ đọc cờ và backoff không có tranh chấp nào để giảm. Cửa sổ quay nhân
đôi sau mỗi lần đọc thấy cờ chưa đảo, nên thread chờ càng lâu thì càng nhận ra cờ đảo muộn.
% nguồn: barrier.csv, ns/episode sr_bo_spin/sr: p=2 2,30; p=8 9,54.

\emph{Độ vững.} Giữa 7 lần lặp, một episode của \PCname{sr} ở 4 thread dao động 163--340~ns. Tính riêng
cho từng lượt lặp, độ dốc từ 4 tới 10 thread vẫn lớn hơn độ dốc từ 2 tới 4 thread ở cả 7 lượt, nhưng gấp
1,2 tới 10,7 lần. Vì vậy hướng tăng của độ dốc là chắc chắn, còn hệ số gần 6 thì không.
% nguồn: barrier.csv sr p=4, 7 lần: 163,0..339,5 ns/episode; độ dốc theo lượt -> tỉ số 1,22..10,7.
